\subsection{The general signed construction and its exact fibres}

\begin{statement}{B.13}{Proved}{Signed normalizers, numerical data, and effective conjugator pairs}
Let $C$ be a profinite group topologically generated by $\rho,\sigma$,
with
\[
 \rho^6=\sigma^4=1,\qquad \omega:=\rho^3=\sigma^2,
\]
and suppose that the automorphism $\tau$ defined on these generators by
\[
 \tau(\sigma)=\sigma^{-1},\qquad
 \tau(\rho)=\sigma\rho^{-1}\sigma^{-1}
                                                        \tag{B.13.1}
\]
exists. Set $e_0=\sigma\rho$, $e_1=\rho\sigma$,
$\ell_1=e_1^2$, and $L_0=\overline{\langle e_0\rangle}$.
Write $A=\widehat\Z^\times$, and let $\epsilon_3:A\to\{\pm1\}$
and $\epsilon_2:A\to\{\pm1\}$ be reduction modulo $3$ and $4$.
The groups of C3:530, (7), and C3:549, (6), are respectively
\[
\begin{aligned}
 B={}&\{(u,\mu): u(e_0)=e_0^\mu,
       \ u(\rho)\sim_C\rho^{\epsilon_3(\mu)},
       \ u(\sigma)\sim_C\sigma^{\epsilon_2(\mu)}\},\\
 Z={}&\{(u,\mu,e_r,e_s):u(e_0)=e_0^\mu,
       \ u(\rho)\sim_C\rho^{e_r},
       \ u(\sigma)\sim_C\sigma^{e_s}\},
\end{aligned}
\tag{B.13.2}
\]
where $u\in\Aut(C)$, $\mu\in A$, and $e_r,e_s\in\{\pm1\}$.
We order the two recorded signs as $(\rho,\sigma)$; the manuscript
often orders them as $(\sigma,\rho)$.
The embedding $B\hookrightarrow Z$ supplies the two displayed
characters of $\mu$. It is normal of index $1$, $2$, or $4$.

For either $D=B$ or $D=Z$, retain these two characters, denoted
$\chi_r,\chi_s$, as well as the multiplier $\chi:D\to A$.
Put
\[
\begin{gathered}
 j(\ell)=(\operatorname{Int}\ell,1,1,1),\quad
 K=L_0\cap Z(C),\quad \gamma_D=D/j(L_0),\\
 G_D=(C\rtimes D)/\{(\ell^{-1},j(\ell)):\ell\in L_0\}.
\end{gathered}
\tag{B.13.3}
\]
In the $B$ case the last two coordinates of $j$ are suppressed.
Then $D\to G_D$ is injective, $\ker(C\to G_D)=K$, and
$1\to C/K\to G_D\xrightarrow{\delta}\gamma_D\to1$ is exact.
Assume $K=1$ from now on. The three characters descend to $\gamma_D$.
The element $t=(\tau,-1,-1,-1)\in D$ has order $2$, and
\[
\begin{gathered}
 C^{\pm}:=C\rtimes\langle t\rangle\hookrightarrow G_D,
 \qquad t'=\sigma^{-1}t,\\
 (t')^2=1,\quad t'\rho(t')^{-1}=\rho^{-1},\quad
 t'\sigma(t')^{-1}=\sigma^{-1}.
\end{gathered}
\tag{B.13.4}
\]
Here $C^{\pm}$ denotes the whole extension.

For $x=\rho,\sigma$, write $C_x=C_C(x)$, $Z_x=C_{G_D}(x)$,
and $\chi_x=\chi_r,\chi_s$, respectively. Define
\[
\begin{aligned}
 E_x&=\{e\in\{\pm1\}:x^e\sim_C x\},\\
 N_x&=\{(g,e)\in G_D\times\{\pm1\}:gxg^{-1}=x^e\},\\
 N_x^*&=\{(g,e)\in N_x:e=\chi_x(\delta g)\},\qquad
 Z_x^*=Z_x\cap\delta^{-1}(\ker\chi_x).
\end{aligned}
\tag{B.13.5}
\]
These formulas realize the external semidirect products $N_x$ of
C3:540--541 and 555 without conflating them with their images in $G_D$.
The precise image and kernel statements are
\[
\begin{gathered}
 \delta(Z_x)=\chi_x^{-1}(E_x),\qquad
 \operatorname{im}(N_x\to G_D)=N_{G_D}(\langle x\rangle),\\
 \ker(N_x\to G_D)=\{(1,e):x^e=x\},
 \qquad N_x^*\longrightarrow G_D\text{ is always injective},\\
 1\longrightarrow C_x\longrightarrow N_x^*
       \xrightarrow{\delta}\gamma_D\longrightarrow1,\\
 1\longrightarrow Z_x^*\longrightarrow N_x^*
       \xrightarrow{e}\{\pm1\}\longrightarrow1.
\end{gathered}
\tag{B.13.6}
\]
The second sequence splits by $-1\mapsto(t',-1)$.
Thus $N_x\to G_D$ has a kernel of order $2$ exactly when $x^2=1$;
the starred map remains injective in that case.

The signed groups of C3:559--560, (50)--(59), are the fibre products
\[
\begin{aligned}
 \Gamma_D&=N_\rho^*\times_{\gamma_D}N_\sigma^*,&
 S_0\Gamma_D&=D\times_{\gamma_D}\Gamma_D,\\
 \mathcal M_D&=\Gamma_D\times_{\gamma_D}G_D,&
 S_0\mathcal M_D&=D\times_{\gamma_D}\mathcal M_D.
\end{aligned}
\tag{B.13.7}
\]
Their kernels over $\gamma_D$ are, respectively,
\[
\begin{gathered}
 C_\rho\times C_\sigma,\quad
 L_0\times C_\rho\times C_\sigma,\\
 C_\rho\times C_\sigma\times C,\quad
 L_0\times C_\rho\times C_\sigma\times C.
\end{gathered}
\tag{B.13.8}
\]
Every arrow to $\gamma_D$ is onto. Forgetting the $D$ coordinate
gives extensions by $L_0$; forgetting the $G_D$ coordinate gives
extensions by $C$. The square
$S_0\mathcal M_D=\mathcal M_D\times_{\Gamma_D}S_0\Gamma_D$
is cartesian.

Finally, put
\[
 \Lambda=\{\mu\in A:e_0^\mu=e_0\}.
\]
The kernels of forgetting numerical coordinates and retaining only
the automorphism $u$ are exactly
\[
\begin{aligned}
 \ker(Z\to\Aut C)&=\Lambda\times E_\rho\times E_\sigma,\\
 \ker(B\to\Aut C)&=
 \{\mu\in\Lambda:\epsilon_3(\mu)\in E_\rho,
                         \epsilon_2(\mu)\in E_\sigma\}.
\end{aligned}
\tag{B.13.9}
\]
They are central. These kernels must not be confused with the
manuscript's distinction between numerical solutions and effective
solutions: ``effective'' there means that the two proposed generator
images extend to an automorphism, and the numerical data are still kept.
There is an exact conjugator parametrization of that effective subset,
described in the proof, with fibres $C_\rho\times C_\sigma$.
\end{statement}
\begin{proof}
The element $\omega$ commutes with both generators and hence with $C$.
Formula (B.13.1) gives $\tau^2=1$ on both generators and gives
$\tau(e_0)=e_0^{-1}$. It also gives $\tau(e_1)=e_1^{-1}$, using
$\sigma^2=\omega$ and centrality of $\omega$.
Composition preserves (B.13.2): for example, if $v(\rho)$ is conjugate
to $\rho^f$ and $u(\rho)$ to $\rho^e$, then $uv(\rho)$ is conjugate
to $\rho^{ef}$. The same calculation for $\sigma$, and
$u(e_0^\nu)=e_0^{\mu\nu}$, proves closure under composition;
applying it to inverses proves the subgroup assertion.
The kernel of
\[
 Z\longrightarrow\{\pm1\}^2,
 \qquad (u,\mu,e_r,e_s)\longmapsto
       (e_r\epsilon_3(\mu)^{-1},e_s\epsilon_2(\mu)^{-1})
\]
is precisely $B$. This proves normality and the possible indices.
The conditions are closed in the product of the profinite groups
$\Aut(C)$ and the numerical factors: conjugacy classes are compact,
and the power map on $L_0$ is continuous. The profiniteness of
$\Aut(C)$ for topologically finitely generated $C$ was justified in
B.11. All groups and arrows in this proof carry these topologies.

Every $u$ in $D$ carries $\ell\in L_0$ to $\ell^{\chi(u)}$.
Consequently $j(L_0)$ is normal in $D$, with kernel $K$ in $L_0$.
In $C\rtimes D$, the antidiagonal in (B.13.3) is an injective image
of $L_0$; it centralizes $C$, and conjugation by $d\in D$ carries
its element indexed by $\ell$ to the one indexed by $u(\ell)$.
It is therefore closed and normal. Intersecting it with $C$ and
with $D$ gives the kernel and injectivity assertions in (B.13.3).
Projection to $D/j(L_0)$ gives the displayed extension. This also
proves the universal property: homomorphisms from $C,D$ respecting
the given action factor through $G_D$ exactly when they identify
$\ell\in C$ with $j(\ell)\in D$ for every $\ell\in L_0$.

The element $t$ belongs to $D$ by (B.13.1), and $t\notin C$ in
$G_D$ because $\chi(t)=-1$, whereas $\chi(C)=1$.
The resulting map $C\rtimes\langle t\rangle\to G_D$ is injective:
its intersection with $C$ is injective, and its two cosets have
distinct multiplier. Direct multiplication gives
$(\sigma^{-1}t)^2=\sigma^{-1}\tau(\sigma^{-1})=1$ and both
inversion identities in (B.13.4).

An element $g\in G_D$ above $d\in\gamma_D$ sends $x$ to a
$C$-conjugate of $x^{\chi_x(d)}$. Some $C$-multiple of $g$
centralizes $x$ if and only if $\chi_x(d)\in E_x$.
This proves the first image formula in (B.13.6), including both
directions. The only generators of a cyclic group of order dividing
$6$ or $4$ are $x,x^{-1}$, with possible equality; this proves the
normalizer image assertion. The kernel formula follows by putting
$g=1$. The sign map on $N_x$ is split by $(t',-1)$, and its kernel
is $Z_x$. Thus this definition agrees with the external semidirect
product $Z_x\rtimes C_2$ even when $t'\in Z_x$.

For a fixed $d\in\gamma_D$, choose $g$ above it and $c\in C$ with
$gxg^{-1}=cx^{\chi_x(d)}c^{-1}$. Then
$(c^{-1}g,\chi_x(d))\in N_x^*$ maps to $d$.
The fibre at $1$ consists exactly of $(c,1)$ with $c\in C_x$.
If $(1,e)\in N_x^*$, its sign is forced to be $\chi_x(1)=1$;
this proves injectivity for every order of $x$. Finally the kernel
of its sign map is $Z_x^*$, and $(t',-1)$ is a splitting because
$\chi_x(\delta t')=-1$. Taking fibre products now proves every
surjectivity, kernel, extension and cartesian assertion of
(B.13.7)--(B.13.8). In particular the upper $Z_x$ in C3:556,
(41), (43), has to mean $Z_x^*$; otherwise it need not be a
subgroup of $N_x^*$. The numerical kernel formulas (B.13.9) follow
by setting $u=1$ in (B.13.2). They are central because the numerical
factors are abelian and their automorphism coordinate is the identity.
No surjectivity of $\chi:\gamma_D\to\widehat\Z^\times$ has been
used or asserted.

We give the promised exact parametrization. For $D=Z$, let $\Sigma$
be the set of $(\xi,\eta,\mu,e_r,e_s)$ such that
\[
 \xi\rho\sim_C\rho^{e_r},\quad
 \eta\sigma\sim_C\sigma^{e_s},\quad
 \xi=\eta\ell_1^{\nu},\qquad \nu=(\mu-1)/2.
                                                        \tag{B.13.10}
\]
For $D=B$, prescribe $e_r=\epsilon_3(\mu)$,
$e_s=\epsilon_2(\mu)$. Division by $2$ here is unique in
$\widehat\Z$ because every profinite unit is odd.
Let $\Sigma^{\mathrm{eff}}$ be the subset for which an automorphism
$u$ sends $(\rho,\sigma)$ to $(\xi\rho,\eta\sigma)$.
The automorphism is unique by topological generation. It fixes
$\omega$, since its image of $\sigma$ is conjugate to
$\sigma^{\pm1}$. For $e'_0=\sigma^{-1}\rho=\omega e_0$, the
equation $u(e_0)=e_0^\mu$ is consequently equivalent to
$u(e'_0)=(e'_0)^\mu$. Multiplying this latter equation gives
\[
 \eta^{-1}\xi=(\rho\sigma^{-1})^{\mu-1}
             =(\rho\sigma)^{\mu-1}=\ell_1^\nu.
\]
The middle equality uses the evenness of $\mu-1$ and $\omega^2=1$.
Thus the map
\[
 d\longmapsto\bigl(u(\rho)\rho^{-1},u(\sigma)\sigma^{-1},
 \mu,e_r,e_s\bigr)
\]
is a bijection $D\to\Sigma^{\mathrm{eff}}$.
This statement keeps the existence condition; it does not identify
all of $\Sigma$ with the effective subset.

Let $P_D^{\mathrm{eff}}$ retain, in addition, conjugators
$\alpha,\beta\in C^{\pm}$ of orientations $e_r,e_s$ with
$u(\rho)=\alpha\rho\alpha^{-1}$ and
$u(\sigma)=\beta\sigma\beta^{-1}$.
Each fibre over $d$ is a nonempty right torsor under
$C_\rho\times C_\sigma$: the class conditions supply conjugators
of the required orientation using $t'$, and two such conjugators
of the same orientation differ by the corresponding centralizer
in $C$. Write $e_r=(-1)^{v_r}$, $e_s=(-1)^{v_s}$.
Using $i_x:N_x^*\hookrightarrow G_D$, the bijection is
\[
\begin{aligned}
 S_0\Gamma_D&\longrightarrow P_D^{\mathrm{eff}},\\
 (d,a,b)&\longmapsto(\alpha,\beta,\mu,e_r,e_s),\\
 \alpha&=d\,i_\rho(a)^{-1}(t')^{v_r},\qquad
 \beta=d\,i_\sigma(b)^{-1}(t')^{v_s}.
\end{aligned}
\tag{B.13.11}
\]
Equal images in $\gamma_D$ show that $d\,i_\rho(a)^{-1}$ and
$d\,i_\sigma(b)^{-1}$ belong to $C$. Inverting the prescribed
actions on $\rho,\sigma$ verifies both conjugation identities.
Conversely, recover $d$ from the effective images and numerical
data, and recover $a,b$ by
\[
 i_\rho(a)=(t')^{v_r}\alpha^{-1}d,\qquad
 i_\sigma(b)=(t')^{v_s}\beta^{-1}d.
                                                        \tag{B.13.12}
\]
These elements have the required signs and the same $\gamma_D$
image as $d$. Their lifts to the starred groups are unique even
if $\rho^2=1$ or $\sigma^2=1$. This proves both inverseness and
the exact fibre assertion. These bijections are homeomorphisms:
the source fibre products are compact, the parameter spaces are
Hausdorff subspaces of the displayed products, and the formulas
are continuous. The parametrization also supplies an associative group
law on the effective conjugator parameters by transporting the
componentwise law on $S_0\Gamma_D$. It does not impose a group law
on numerical solutions that fail realization.

These are statements about profinite groups. For the finite shadows
in B.14 they give, after taking constant finite group schemes over
$\Z$, exact sequences of fppf sheaves: every quotient map of finite
constant groups is already locally split on each constant component.
The same formal kernel arguments apply to group objects in a sheaf
category with locally stated conjugacy hypotheses and the indicated
quotients. This does not prove representability or finite presentation
for an arbitrary group-scheme analogue of $C$; nor does it identify
the full noncongruence groups with the matrix schemes of \S48.
\end{proof}

\begin{statement}{B.14}{Refuted}{Unrestricted centralizers and a left sign factor do not give the signed formulas}
The unqualified inclusions $Z_x\subset N_x^*$ implicit in the upper
rows of C3:556, (41), (43), are false in general. The exact replacement
is $Z_x^*$ in (B.13.5). Also, with conjugation $g(x)=gxg^{-1}$,
the change of conjugator in C3:546, fol.~576, (64), must be
\[
                 \beta'=\beta(t')^v,\qquad e_s=(-1)^v,
                                                        \tag{B.14.1}
\]
to obtain
$\operatorname{Int}(\beta')(\sigma)
 =\omega^v\operatorname{Int}(\beta)(\sigma)$.
The printed left factor $(t')^v\beta$ does not have that property.

The following are exact finite checks of the corrected signed
construction. We replace $\widehat\Z^\times$ by
$U_{12}=(\Z/12)^\times$ throughout the table, take
$\rho=(123),\sigma=(12)$ in $S_3$, and
$\rho=(123),\sigma=(34)$ in $S_4$. In the two matrix groups take
\[
 \rho=\begin{pmatrix}0&1\\-1&1\end{pmatrix},\qquad
 \sigma=\begin{pmatrix}0&-1\\1&0\end{pmatrix}.
\]
The symbols in each row refer to $D=B_{12}$ or $Z_{12}$,
respectively, and to (B.13.3)--(B.13.8):
\begingroup\small\setlength{\arraycolsep}{4pt}
\[
\begin{array}{c|c|rrrrrrr}
 C&D&|D|&|\gamma_D|&|G_D|&|N_\rho^*|&|N_\sigma^*|
       &|\mathcal M_D|&|S_0\mathcal M_D|\\\hline
 S_3&B&8&4&24&12&8&144&288\\
 S_3&Z&32&16&96&48&32&576&1152\\
 S_4&B&16&4&96&12&16&1152&4608\\
 S_4&Z&64&16&384&48&64&4608&18432\\
 \mathrm{SL}_2(\F_3)&B&12&4&96&24&16&2304&6912\\
 \mathrm{SL}_2(\F_3)&Z&24&8&192&48&32&4608&13824\\
 \mathrm{SL}_2(\Z/4)&B&16&4&192&24&32&9216&36864\\
 \mathrm{SL}_2(\Z/4)&Z&32&8&384&48&64&18432&73728
\end{array}
\]
\endgroup
In all eight rows $\Sigma=\Sigma^{\mathrm{eff}}$, every fibre
in (B.13.6) has the stated order, and (B.13.11) is bijective.
These equalities of solution sets are finite-case results.
\end{statement}
\begin{proof}
First take $C=S_3$ as in the table. Here $\tau=1$ on $C$:
$\tau$ fixes $\sigma$, and
$\sigma\rho^{-1}\sigma^{-1}=\rho$.
Thus $t=(1,-1)\in B$ centralizes all of $C$, but
$\chi_r(\delta t)=\chi_s(\delta t)=-1$.
Its positive-sign lift $(t,+1)\in N_x$ is not in $N_x^*$.
This gives an explicit counterexample for both $x=\rho,\sigma$,
also in the finite $U_{12}$ shadow. For $\sigma^2=1$, the map
$N_\sigma\to G_D$ has its additional kernel of order $2$;
the starred map has no such kernel. Thus neither a blind use of
centralizers nor a blind identification of the external normalizer
with its image is valid. Formula (B.13.6) is the proved correction.

For (B.14.1), multiplication on the right gives
\[
 \beta(t')^v\sigma(t')^{-v}\beta^{-1}
    =\beta\sigma^{(-1)^v}\beta^{-1}
    =\omega^v\beta\sigma\beta^{-1},
\]
using $\sigma^{-1}=\omega\sigma$. The left factor instead
conjugates the entire element $\beta\sigma\beta^{-1}$ by $t'$.
The distinction is witnessed over $\F_3$ by $\beta=\rho$ and
$t'=\left(\begin{smallmatrix}0&1\\1&0\end{smallmatrix}\right)$:
\[
 (t'\beta)\sigma(t'\beta)^{-1}
       =\begin{pmatrix}2&2\\2&1\end{pmatrix},\qquad
 \omega\beta\sigma\beta^{-1}
       =(\beta t')\sigma(\beta t')^{-1}
       =\begin{pmatrix}2&1\\1&1\end{pmatrix}.
\]
The attached photograph of scan 546 was inspected: its (64)
places the factor $\tau'^v$ on the left. The calculation diagnoses
a mathematical slip under the stated conjugation convention;
it does not claim a final editorial reading of overwritten text.
The inverse in C3:544, (54), (56), must likewise be kept as
$\beta=u\,i_\sigma(b)^{-1}$ if the equality of $\delta$-images
in its (57) is retained. The $b\in Z_\sigma$ printed under
C3:545, (60), must be $b\in N_\sigma$, as its sign condition
otherwise forces $\epsilon_2(\mu)=1$.

We describe the executed finite verification so that its scope is
reproducible. The standalone program \texttt{AB\_signed\_finite.py}
builds the full multiplication tables of the four groups. For every
possible pair of generator images of the correct orders, it
propagates the proposed map through every edge of the right Cayley
graph, rejects an inconsistent assignment, and retains exactly the
bijective assignments. It obtains respectively $6,24,24,48$
automorphisms. It constructs $B_{12},Z_{12}$ from their definitions,
checks the antidiagonal normal subgroup, and forms the actual
semidirect-product quotient $G_D$. It enumerates $Z_x,N_x,N_x^*$,
checks every $\gamma_D$ fibre and every kernel above, and compares
all values and inverse values of (B.13.11)--(B.13.12). It also
enumerates all numerical solutions of (B.13.10), so the asserted
effectivity in these cases is checked independently of the
fibre-product parametrization. Its JSON contains the additional
orders of $\Gamma_D,S_0\Gamma_D$, unrestricted normalizers,
centralizers, and numerical kernels. The largest displayed product
has $73{,}728$ elements by these exact fibre counts.

For clarity, the inverse-conjugacy pairs $(E_\rho,E_\sigma)$ are
\[
 (\{\pm1\},\{\pm1\}),\quad
 (\{\pm1\},\{\pm1\}),\quad
 (\{1\},\{\pm1\}),\quad
 (\{\pm1\},\{1\})
\]
in the order of the four groups. The numerical kernel orders of
$B_{12}\to\Aut(C)$ are $4,2,2,2$; those of
$Z_{12}\to\Aut(C)$ are $16,8,4,4$. These calculations verify that
forgetting the numerical parameters is not injective even though
every numerical solution is effective in these examples.

All powers in the definitions factor through $U_{12}$ because
$|e_0|=2,4,3,4$, respectively. The full profinite objects are
$B=B_{12}\times_{U_{12}}\widehat\Z^\times$ and
$Z=Z_{12}\times_{U_{12}}\widehat\Z^\times$, not the finite groups
in the table. Reduction $\widehat\Z^\times\to U_{12}$ is onto:
choose the required units at $2$ and $3$ and the unit $1$ at every
other prime. The finite table therefore verifies every signed
condition relevant to these quotients without confusing a finite
numerical shadow with the full profinite construction.
\end{proof}
